\begin{frame}{한 편의 MDP를 통째로 쓰기: 3칸 보드게임}
  추상적으로 다섯 요소를 나열하는 것만으로는 체감되지 않는다.
  아주 작은 MDP 하나를 \textbf{통째로} 써본다.
  \vskip0.4em
  \kb{설정}: 상태 0, 1, 2의 보드게임. 0에서 시작, 매 스텝
  ``오른쪽으로 한 칸'' 또는 ``자리참기''를 고름, 상태 2에 도착하면
  \textbf{에피소드 종료}. 이동은 확률적: ``오른쪽''을 골라도 0.3의
  확률로 제자리에 미끄러진다.
  \vskip0.4em
  \begin{itemize}
    \item $\mathcal{S} = \{0, 1, 2\}$ --- 보드의 세 칸. 2는
      \kb{터미널 상태}(terminal state, 게임 종료 상태).
    \item $\mathcal{A} = \{\text{right, wait}\}$ --- 두 행동.
    \item $\gamma = 0.9$.
  \end{itemize}
  \vskip0.4em
  {\footnotesize 이 다섯 개(5-tuple)가 쓰이는 순간, 이 게임의 ``규칙''은
  \textbf{완전히} 정해졌다. 어떤 시작점에서든, 어떤 정책(상태마다
  행동을 고르는 규칙)으로 놀든, 받을 보상 시퀀스의 \textbf{전체
  확률분포}가 이 다섯 요소만으로 계산된다 --- 게임에 ``보이지 않는
  상태''가 없어야 한다는 전제(= 마르코프 성질)를 붙이면 정확하다.
  다음 장에서 벨만방정식으로 $V^\pi(s)$ 를 계산하면, 여기 쓴 $P$ 와
  $R$ 이 그대로 수식에 대입된다.}
\end{frame}
